\documentclass[12pt, a4paper, oneside]{ctexart}
\usepackage{amsmath, amsthm, amssymb, bm, xcolor, tikz, framed, graphicx, hyperref, mathrsfs}

\title{Soild State Physics Homework (Week 1-2)}
\author{Tang Yi \\ [学号已脱敏]}
\usepackage[a4paper, left=2cm, right=2cm, top=2cm, bottom=2cm]{geometry}
\linespread{1.3} % 行距更紧凑
\definecolor{shadecolor}{RGB}{241, 241, 255}

\newcounter{problemname}
\newenvironment{problem}
  {\begin{shaded}\stepcounter{problemname}\par\noindent\textbf{题目 \theproblemname. }}
  {\end{shaded}\par}
\newenvironment{solution}
  {\par\noindent\textbf{题目 \theproblemname 的解答. }}
  {\par}
\newenvironment{note}
  {\par\noindent\textbf{题目 \theproblemname 的注记. }}
  {\par}

\begin{document}

\maketitle

\begin{problem}
    如果将等体积球分别排成下列结构，设 $x$ 表示刚球所占体积与总体积之比，证明
    
    \begin{center}
    \begin{tabular}{cc}
        结构 & $x$ \\
        简单立方 & $\pi/6 \approx 0.52$ \\
        体心立方 & $\sqrt{3}\pi/8 \approx 0.68$ \\
        面心立方 & $\sqrt{2}\pi/6 \approx 0.74$ \\
        六方密排 & $\sqrt{2}\pi/6 \approx 0.74$ \\
        金刚石 & $\sqrt{3}\pi/16 \approx 0.34$
    \end{tabular}
    \end{center}
\end{problem}

\begin{solution}

设原子为半径 $R$ 的刚球，单位晶胞体积为 $V_c$，晶胞内原子数为 $N$，则堆积因子为

\[
x=\frac{N \cdot \frac{4}{3}\pi R^3}{V_c}.
\]

\textbf{(1) 简单立方结构 (SC)}

简单立方晶格中 8 个顶点原子各贡献 $1/8$，故

\[
N=1.
\]

原子沿立方体棱方向接触，因此

\[
a=2R.
\]

晶胞体积

\[
V_c=a^3=(2R)^3=8R^3.
\]

因此

\[
x=\frac{\frac{4}{3}\pi R^3}{8R^3}
=\frac{\pi}{6}\approx0.52.
\]

\textbf{(2) 体心立方结构 (BCC)}

单位胞内原子数

\[
N=2.
\]

原子沿体对角线接触：

\[
\sqrt{3}a=4R.
\]

因此

\[
a=\frac{4R}{\sqrt3}.
\]

晶胞体积

\[
V_c=a^3=\left(\frac{4R}{\sqrt3}\right)^3
=\frac{64R^3}{3\sqrt3}.
\]

堆积因子

\[
x=\frac{2\cdot\frac{4}{3}\pi R^3}{V_c}
=\frac{\sqrt3\pi}{8}
\approx0.68.
\]

\textbf{(3) 面心立方结构 (FCC)}

单位胞原子数

\[
N=4.
\]

原子沿面对角线接触

\[
\sqrt2 a=4R.
\]

因此

\[
a=2\sqrt2 R.
\]

晶胞体积

\[
V_c=a^3=(2\sqrt2R)^3=16\sqrt2R^3.
\]

堆积因子

\[
x=\frac{4\cdot\frac{4}{3}\pi R^3}{16\sqrt2R^3}
=\frac{\sqrt2\pi}{6}
\approx0.74.
\]

\textbf{(4) 六方密排结构 (HCP)}

HCP 与 FCC 具有相同的堆积效率，因此

\[
x=\frac{\sqrt2\pi}{6}\approx0.74.
\]

\textbf{(5) 金刚石结构}

金刚石结构为 FCC 晶格加两原子基元，单位胞中共有

\[
N=8
\]

个原子。

最近邻距离为

\[
d=\frac{\sqrt3}{4}a.
\]

由于原子接触

\[
2R=d,
\]

故

\[
R=\frac{\sqrt3}{8}a.
\]

单位胞体积

\[
V_c=a^3.
\]

球体总体积为

\[
8\cdot\frac{4}{3}\pi R^3.
\]

因此

\[
x=\frac{8\cdot\frac{4}{3}\pi R^3}{a^3}
=\frac{\sqrt3\pi}{16}
\approx0.34.
\]

\end{solution}

\begin{problem}
    试证六方密排密堆积结构中 $\dfrac{c}{a} = \left( \dfrac{8}{3} \right)^{1/2} \approx 1.633$.
\end{problem}

\begin{solution}

考虑六方密排结构中三层原子堆积。底层三个原子中心形成边长为 $2R$ 的正三角形。

设上层原子中心到该三角形平面的高度为 $h$。三角形中心到顶点距离为

\[
r=\frac{2R}{\sqrt3}.
\]

由于上层原子与三个底层原子相接触，因此球心距离满足

\[
(2R)^2=h^2+r^2.
\]

代入 $r$ 得

\[
4R^2=h^2+\frac{4R^2}{3}.
\]

于是

\[
h^2=\frac{8R^2}{3},
\qquad
h=\sqrt{\frac{8}{3}}R.
\]

HCP 单胞高度为

\[
c=2h,
\]

而

\[
a=2R.
\]

因此

\[
\frac{c}{a}
=\frac{2h}{2R}
=\frac{h}{R}
=\sqrt{\frac{8}{3}}
\approx1.633.
\]

\end{solution}

\begin{problem}
    证明：体心立方晶格的倒格子是面心立方；面心立方晶格的倒格子是体心立方.
\end{problem}

\begin{solution}

倒格子基矢定义为

\[
\boldsymbol b_1
=2\pi
\frac{\boldsymbol a_2\times\boldsymbol a_3}
{\boldsymbol a_1\cdot(\boldsymbol a_2\times\boldsymbol a_3)},
\]

\[
\boldsymbol b_2
=2\pi
\frac{\boldsymbol a_3\times\boldsymbol a_1}
{\boldsymbol a_1\cdot(\boldsymbol a_2\times\boldsymbol a_3)},
\]

\[
\boldsymbol b_3
=2\pi
\frac{\boldsymbol a_1\times\boldsymbol a_2}
{\boldsymbol a_1\cdot(\boldsymbol a_2\times\boldsymbol a_3)}.
\]

\textbf{(1) BCC 的倒格子}

取 BCC 原胞基矢

\[
\boldsymbol a_1=\frac a2(1,1,-1),\qquad
\boldsymbol a_2=\frac a2(1,-1,1),\qquad
\boldsymbol a_3=\frac a2(-1,1,1).
\]

计算得到

\[
\boldsymbol b_1=\frac{2\pi}{a}(0,1,1),\quad
\boldsymbol b_2=\frac{2\pi}{a}(1,0,1),\quad
\boldsymbol b_3=\frac{2\pi}{a}(1,1,0).
\]

这正是面心立方晶格的原胞基矢，因此

\[
\text{BCC 的倒格子为 FCC}.
\]

\textbf{(2) FCC 的倒格子}

FCC 原胞基矢为

\[
\boldsymbol a_1=\frac a2(0,1,1),\qquad
\boldsymbol a_2=\frac a2(1,0,1),\qquad
\boldsymbol a_3=\frac a2(1,1,0).
\]

计算得到

\[
\boldsymbol b_1=\frac{2\pi}{a}(1,-1,1),\quad
\boldsymbol b_2=\frac{2\pi}{a}(1,1,-1),\quad
\boldsymbol b_3=\frac{2\pi}{a}(-1,1,1).
\]

这正是体心立方晶格的原胞基矢，因此

\[
\text{FCC 的倒格子为 BCC}.
\]

\end{solution}

\begin{problem}
    证明：倒格子原胞的体积为 $(2\pi)^3/v_c$，其中 $v_c$ 为正格子原胞的体积.
\end{problem}

\begin{solution}

正格子原胞体积为

\[
V_c=\boldsymbol a_1\cdot(\boldsymbol a_2\times\boldsymbol a_3).
\]

倒格子基矢

\[
\boldsymbol b_1
=2\pi\frac{\boldsymbol a_2\times\boldsymbol a_3}{V_c},\quad
\boldsymbol b_2
=2\pi\frac{\boldsymbol a_3\times\boldsymbol a_1}{V_c},\quad
\boldsymbol b_3
=2\pi\frac{\boldsymbol a_1\times\boldsymbol a_2}{V_c}.
\]

倒格子原胞体积

\[
V^*=\boldsymbol b_1\cdot(\boldsymbol b_2\times\boldsymbol b_3).
\]

代入上式并整理可得

\[
V^*=(2\pi)^3
\frac{(\boldsymbol a_2\times\boldsymbol a_3)\cdot
[(\boldsymbol a_3\times\boldsymbol a_1)\times(\boldsymbol a_1\times\boldsymbol a_2)]}
{V_c^3}.
\]

利用向量恒等式

\[
(\boldsymbol a_2\times\boldsymbol a_3)\cdot
[(\boldsymbol a_3\times\boldsymbol a_1)\times(\boldsymbol a_1\times\boldsymbol a_2)]
=V_c^2
\]

最终得到

\[
V^*=\frac{(2\pi)^3}{V_c}.
\]

\end{solution}

\begin{problem}
    证明：倒格子矢量 $\boldsymbol{G} = h_1 \boldsymbol{b}_1 + h_2 \boldsymbol{b}_2 + h_3 \boldsymbol{b}_3$ 垂直于密勒指数为 $(h_1 h_2 h_3)$ 的晶面系.
\end{problem}

\begin{solution}

设倒格子矢量

\[
\boldsymbol G=h_1\boldsymbol b_1+h_2\boldsymbol b_2+h_3\boldsymbol b_3.
\]

具有密勒指数 $(h_1h_2h_3)$ 的晶面与晶轴的截距为

\[
\frac{a_1}{h_1},\qquad
\frac{a_2}{h_2},\qquad
\frac{a_3}{h_3}.
\]

该晶面方程可以写为

\[
\boldsymbol G\cdot\boldsymbol r=2\pi.
\]

而一般平面方程为

\[
\boldsymbol n\cdot\boldsymbol r=\text{const},
\]

其中 $\boldsymbol n$ 为平面的法向量。

比较两式可知

\[
\boldsymbol n=\boldsymbol G.
\]

因此倒格子矢量

\[
\boldsymbol G=h_1\boldsymbol b_1+h_2\boldsymbol b_2+h_3\boldsymbol b_3
\]

垂直于密勒指数为 $(h_1h_2h_3)$ 的晶面系。

\end{solution}

\end{document}
